\documentclass[11pt, a4paper]{article}

% Gemeinsame Style-Definitionen (TEXINPUTS muss templates/ enthalten — siehe scripts/build_tex.py)
\input{bfw-style.tex}

\title{\vspace*{-1.5cm}\Huge\bfseries\color{bfwPrimaryDark}Session~1: IT-Berufe \& duales System\\[0.4em]
  \large\color{bfwInkSoft}\textnormal{Von der Neuordnung 2020 bis zur gestreckten Abschlussprüfung}}
\author{\textsf{IT Lernfeld 1 --- Das Unternehmen und die eigene Rolle im Betrieb beschreiben}}
\date{}

\begin{document}
\maketitle
\thispagestyle{empty}

\vspace{1em}
\begin{merksatz}[title={\bfseries Worum geht es in dieser Session?}]
Sie beginnen Ihre Ausbildung in einem IT-Beruf und stehen damit zugleich im \emph{dualen
System der Berufsausbildung}. Dieses Handout klärt, welche vier IT-Ausbildungsberufe es
seit der Neuordnung 2020 gibt, was den Fachinformatiker mit seinen vier Fachrichtungen
auszeichnet und wie das duale System mit seinen zwei Lernorten funktioniert. Sie lernen die
Beteiligten --- Betrieb, Berufsschule und IHK --- und ihre Aufgaben kennen sowie die
Struktur der gestreckten Abschlussprüfung. Das ist die Grundlage für Lernfeld~1 und für
die IHK-Prüfung.
\end{merksatz}

\tableofcontents
\vspace{1em}
\hrule
\vspace{1em}

\section{Die vier IT-Ausbildungsberufe seit 2020}

Seit 1969 wird in Berufen der Datenverarbeitung ausgebildet. Mit der \textbf{Neuordnung
2020} wurden die IT-Berufe grundlegend reformiert. Seitdem gibt es vier anerkannte
IT-Ausbildungsberufe, die gemeinsam die \emph{Berufsgruppe IT-Berufe} bilden:

\begin{enumerate}
  \item \textbf{IT-Systemelektroniker/-in}
  \item \textbf{Kaufmann/-frau für IT-Systemmanagement}
  \item \textbf{Kaufmann/-frau für Digitalisierungsmanagement}
  \item \textbf{Fachinformatiker/-in} (mit vier Fachrichtungen)
\end{enumerate}

Alle vier Berufe teilen in den ersten eineinhalb Jahren gemeinsame
\emph{Kernqualifikationen (KQ)} --- berufsübergreifende Inhalte aus den Lernfeldern
1 bis 6. Das ermöglicht eine gemeinsame Beschulung und bereitet auf Teil~1 der
gestreckten Abschlussprüfung vor. Im weiteren Verlauf wachsen die berufsprofilgebenden
\emph{Fachqualifikationen (FQ)}. Im dritten Ausbildungsjahr werden Auszubildende
zunehmend selbstständig in ausgewählten Einsatzgebieten tätig --- dort entsteht auch die
Projektaufgabe für die Abschlussprüfung (vgl. §§~7, 8, 45 BBiG).

\begin{center}
\begin{tabular}{p{5cm} p{8cm}}
\toprule
\textbf{Beruf} & \textbf{Schwerpunkt / typische Tätigkeiten}\\
\midrule
IT-Systemelektroniker/-in & Herstellen und Betreuen von IT-Infrastrukturen; Service und Instandsetzung; elektrische Sicherheit und IT-Sicherheit; Beratung\\
Kaufmann/-frau für IT-Systemmanagement & Vermarktung und Administration von IT-Lösungen; Angebote und Auftragsbearbeitung; Kundenberatung; IT-Beschaffung\\
Kaufmann/-frau für Digitalisierungsmanagement & Daten und Prozesse aus betriebswirtschaftlicher Perspektive; Management der Digitalisierung von Geschäftsprozessen\\
Fachinformatiker/-in & Konzeption, Entwicklung, Installation oder Analyse von IT-Systemen und Software (je nach Fachrichtung)\\
\bottomrule
\end{tabular}
\end{center}

\begin{merksatz}
Der \textbf{Fachinformatiker} ist laut Lehrbuch der am meisten nachgefragte IT-Beruf,
weil ein hoher Bedarf an Systemintegration, Softwareentwicklung und Vernetzung besteht.
\end{merksatz}

\section{Die vier Fachrichtungen des Fachinformatikers}

Wegen der vielfältigen Aufgaben und Einsatzbereiche wird der Fachinformatiker in
\textbf{vier Fachrichtungen} ausgebildet. Zwei Fachrichtungen --- \emph{Digitale
Vernetzung} und \emph{Daten- und Prozessanalyse} --- sind mit der Neuordnung 2020
\emph{neu hinzugekommen}, da hier zukünftig hoher Bedarf erwartet wird:

\subsection{Anwendungsentwicklung (FIAE)}

Fachinformatiker dieser Fachrichtung konzipieren und entwickeln Softwarelösungen für
interne und externe Kunden. Sie beachten IT-Sicherheit und Datenschutz, planen Aufgaben
in Abstimmung mit Geschäftsprozessen, beurteilen marktgängige IT-Systeme und sind auch
am Auftragsabschluss beteiligt.

\subsection{Systemintegration (FISI)}

Sie konzipieren, installieren und administrieren vernetzte IT-Systeme. Sie informieren und
beraten Kunden, programmieren Softwarelösungen, betreuen IT-Systeme und sind für IT-Sicherheit
und Datenschutz zuständig.

\subsection{Digitale Vernetzung (FIDV) --- neu seit 2020}

Sie analysieren und planen Systeme zur Vernetzung und Automation, insbesondere in
Produktionsbetrieben. Sie arbeiten mit Netzwerkinfrastruktur und Schnittstellen zwischen
Netzwerkkomponenten und cyber-physischen Systemen. Sie sichern Daten gegen unerlaubte
Zugriffe und beheben Systemausfälle.

\subsection{Daten- und Prozessanalyse (FIDP) --- neu seit 2020}

Sie analysieren Arbeits- und Geschäftsprozesse sowie Datenquellen und stellen
Datenverfügbarkeit und -qualität sicher. Sie nutzen Daten zur Optimierung von Prozessen
und Geschäftsmodellen und entwickeln IT-Lösungen für digitale Produktions- und Geschäftsprozesse.

\section{Das duale System der Berufsausbildung}

\begin{merksatz}[title={\bfseries Ziel der Berufsbildung nach \S~1 BBiG}]
Die Berufsausbildung hat die für die Ausübung einer qualifizierten beruflichen Tätigkeit
in einer sich wandelnden Arbeitswelt notwendigen beruflichen Fertigkeiten, Kenntnisse und
Fähigkeiten (berufliche Handlungsfähigkeit) in einem geordneten Ausbildungsgang zu
vermitteln. Sie hat ferner den Erwerb der erforderlichen Berufserfahrungen zu ermöglichen.
\end{merksatz}

In Deutschland ist die Berufsausbildung staatlich über das
\textbf{Berufsbildungsgesetz (BBiG)} geregelt. Die \emph{duale Berufsausbildung} findet
gleichzeitig an zwei Lernorten statt und ist in über 300 staatlich anerkannten
Ausbildungsberufen möglich.

\subsection{Lernort Betrieb}

Der Betrieb vermittelt die praktische Ausbildung auf Basis einer \emph{Ausbildungsordnung}
(bundeseinheitlich, nach BBiG) und eines individuellen \emph{Ausbildungsplans}. Für die
Ausbildung müssen Betriebe bei der zuständigen IHK zugelassen sein. Der Ausbildungsvertrag
ist schriftlich zu schließen (§~4 Abs.~1 BBiG) und enthält alle wichtigen Vereinbarungen.

\subsection{Lernort Berufsschule (BBS)}

Die Berufsschule erteilt theoretischen Unterricht in beruflichen Lernfeldern und
allgemeinbildenden Fächern (Deutsch, Englisch, Politik/WISO, Sport). Grundlage ist der von
der KMK festgelegte \emph{Rahmenlehrplan}. Der Berufsschulunterricht kann je nach
Bundesland ein bis zwei Tage pro Woche oder im Blockunterricht stattfinden.

\begin{center}
\begin{tabular}{p{6.5cm} p{6.5cm}}
\toprule
\textbf{Betrieb} & \textbf{Berufsschule (BBS)}\\
\midrule
Ausbildungsordnung (BBiG) & Rahmenlehrplan (KMK)\\
Betrieblicher Ausbildungsplan & Didaktische Jahresplanung\\
Praktische Ausbildung & Fachtheorie + Allgemeinbildung\\
Aufsicht: IHK & Aufsicht: Schulleitung / Schulbehörde\\
\bottomrule
\end{tabular}
\end{center}

\section{Beteiligte der Berufsausbildung}

\begin{center}
\begin{tabular}{p{3.5cm} p{9cm}}
\toprule
\textbf{Beteiligter} & \textbf{Aufgaben}\\
\midrule
Auszubildende:r & Erlernt nach BBiG und Ausbildungsordnung einen anerkannten Ausbildungsberuf; führt schriftlichen Ausbildungsnachweis\\
Ausbildungsbetrieb (Ausbildender) & Stellt Auszubildende ein; schließt Berufsausbildungsvertrag; erstellt Ausbildungsplan; muss bei IHK zugelassen sein\\
Ausbilder:in & Führt die Ausbildung im Betrieb durch; muss persönlich und fachlich geeignet und mindestens 24 Jahre alt sein\\
Berufsschule & Erteilt Unterricht nach Rahmenlehrplan; erstellt Jahresplanung; stellt Berufsschulzeugnis aus\\
IHK & Zuständige Stelle nach BBiG; lässt Betriebe zur Ausbildung zu; nimmt Prüfungen ab; berät Beteiligte\\
\bottomrule
\end{tabular}
\end{center}

\begin{merksatz}
Verwechseln Sie nicht \emph{Ausbildende:r} (= der Betrieb als Vertragspartner) und
\emph{Ausbilder:in} (= die Person, die Sie im Betrieb fachlich anleitet). Der Betrieb ist
Ihr Vertragspartner, der Ausbilder ist Ihr Anleiter.
\end{merksatz}

Nach § 278 BGB trägt der Betrieb im Außenverhältnis für das Verschulden seiner Mitarbeiter
--- und damit auch der Auszubildenden --- die Verantwortung. Kann dem Mitarbeiter jedoch
Fahrlässigkeit oder Vorsatz nachgewiesen werden, kann der Betrieb Schadensersatz fordern.
Deshalb gilt: Insbesondere zu Beginn sollten Aufgaben nur nach Anweisung und unter Aufsicht
bearbeitet werden.

\section{Ausbildungsdauer und gestreckte Abschlussprüfung}

\subsection{Ausbildungsdauer und Verkürzung}

Die Regelausbildungsdauer beträgt \textbf{drei Jahre}. Eine Verkürzung ist auf Antrag bei
Anrechnung anerkannter Bildungsgänge möglich (z.\,B.\ bei Hochschulreife oder
Studienabbruch), eine Verlängerung bei besonderer Begründung ebenfalls. Eine Teilzeitausbildung
ist ebenfalls möglich (§§~7, 8, 45 BBiG).

\subsection{Gestreckte Abschlussprüfung der IHK}

Die Abschlussprüfung ist \emph{gestreckt}: Sie besteht aus zwei Teilen und erstreckt sich
über die gesamte Ausbildungszeit.

\begin{center}
\begin{tabular}{p{3.5cm} p{9cm}}
\toprule
\textbf{Prüfungsteil} & \textbf{Details}\\
\midrule
\textbf{Teil 1 (20\,\%)} & Schriftlich, nach etwa 15 Monaten, 90 Minuten; Leitthema „Einrichten eines IT-gestützten Arbeitsplatzes" (Lernfelder 1--6); für alle IT-Berufe gleich\\
\textbf{Teil 2 (80\,\%)} & Zum Ausbildungsende; enthält:\\
& \quad Projektarbeit mit Fachgespräch (50\,\%)\\
& \quad Zwei berufsspezifische Aufgaben (je 10\,\%)\\
& \quad WISO-Prüfung (10\,\%)\\
\bottomrule
\end{tabular}
\end{center}

Die Prüfung gilt als bestanden, wenn das Gesamtergebnis und Teil~2 mindestens
„ausreichend" sind (ab 50\,\%). Eine \emph{Ergänzungsprüfung} von 15 Minuten ist möglich,
wenn damit ein mit mangelhaft bewerteter Prüfungsbereich noch ausgeglichen werden kann
(Gewichtung bisherige Note zu mündlicher Note: 2:1).

\section{Handlungskompetenz und Selbstständiges Lernen}

Der Wandel der Arbeitswelt erfordert von IT-Fachkräften zunehmend
\textbf{lebenslanges Lernen}. Neue Technologien und Geschäftsmodelle verlangen
Flexibilität und Lernbereitschaft. Neben der Fachkompetenz (ca. 50--60\,\% des Berufserfolgs
laut internationalen Studien) gewinnen \emph{Soft Skills} --- persönliche Eigenschaften und
Kompetenzen --- immer mehr an Bedeutung.

\begin{merksatz}[title={\bfseries Handlungskompetenz (nach KMK)}]
Handlungskompetenz wird als die Bereitschaft und Befähigung des Einzelnen verstanden,
sich in beruflichen, privaten und gesellschaftlichen Situationen sachgerecht, individuell,
durchdacht und sozial verantwortlich zu verhalten. Sie entfaltet sich in den Dimensionen
\textbf{Fach-, Sozial- und Selbstkompetenz}. Methodenkompetenz, kommunikative Kompetenz
und Lernkompetenz sind immanenter Bestandteil aller drei Dimensionen.
\end{merksatz}

\begin{center}
\begin{tabular}{p{4.5cm} p{4.5cm} p{4.5cm}}
\toprule
\textbf{Fachkompetenz} & \textbf{Sozialkompetenz} & \textbf{Selbstkompetenz}\\
\midrule
Fachkenntnisse & Kooperations-/Teamfähigkeit & Selbstständigkeit\\
Fertigkeiten & Höflichkeit, Freundlichkeit & Zuverlässigkeit\\
Arbeitsleistung & Konfliktfähigkeit & Sorgfalt\\
Arbeitsqualität & Toleranz & Flexibilität\\
\bottomrule
\end{tabular}
\end{center}

\section{Zusammenfassung}

\begin{enumerate}
  \item Seit \textbf{2020} gibt es vier neu geordnete IT-Berufe; der Fachinformatiker wird in \textbf{vier Fachrichtungen} ausgebildet (zwei davon neu: Digitale Vernetzung, Daten- und Prozessanalyse).
  \item Alle IT-Berufe teilen Kernqualifikationen in \textbf{Lernfeldern 1--6}; ab Lernfeld~7 wird der Unterricht berufsspezifisch.
  \item Das \textbf{duale System} bildet an zwei Lernorten aus: Betrieb (Praxis) und Berufsschule (Theorie); Rechtsgrundlage ist das \textbf{BBiG}.
  \item Beteiligte sind \textbf{Auszubildende}, \textbf{Ausbildungsbetrieb/Ausbilder}, \textbf{Berufsschule} und \textbf{IHK} --- mit klar getrennten Aufgaben.
  \item Ausbildungsdauer: \textbf{3 Jahre}; Verkürzung und Verlängerung nach \S\S~7, 8, 45 BBiG möglich.
  \item Die \textbf{gestreckte Abschlussprüfung} besteht aus Teil~1 (20\,\%, nach 15 Monaten) und Teil~2 (80\,\%, zum Ausbildungsende mit Projektarbeit).
  \item \textbf{Handlungskompetenz} --- Fach-, Sozial- und Selbstkompetenz --- ist das übergeordnete Ausbildungsziel des dualen Systems.
\end{enumerate}

\newpage

\section*{Glossar}
\addcontentsline{toc}{section}{Glossar}

\glossareintrag{Berufsbildungsgesetz (BBiG)}{Gesetz, das die Berufsausbildung in Deutschland regelt; Rechtsgrundlage des dualen Systems.}

\glossareintrag{Duales System}{Berufsausbildung an zwei Lernorten gleichzeitig: Betrieb (Praxis) und Berufsschule (Theorie).}

\glossareintrag{Ausbildende:r}{Der Betrieb als Vertragspartner; stellt Auszubildende ein und schließt den Berufsausbildungsvertrag ab.}

\glossareintrag{Ausbilder:in}{Person im Betrieb, die die Ausbildung praktisch durchführt; muss persönlich und fachlich geeignet und mindestens 24 Jahre alt sein.}

\glossareintrag{Kernqualifikation (KQ)}{Berufsübergreifende Ausbildungsinhalte, die alle vier IT-Berufe in Lernfeld 1 bis 6 gemeinsam erwerben.}

\glossareintrag{Fachqualifikation (FQ)}{Berufsprofilgebende Inhalte, die den spezifischen Beruf von den anderen unterscheiden.}

\glossareintrag{Ausbildungsordnung}{Für jeden anerkannten Ausbildungsberuf nach BBiG erlassene Verordnung; enthält Berufsbild, Dauer und Ausbildungsrahmenplan.}

\glossareintrag{Rahmenlehrplan}{Von der KMK festgelegter Plan, der den schulischen Unterricht in beruflichen Lernfeldern beschreibt (www.kmk.org).}

\glossareintrag{Gestreckte Abschlussprüfung}{IHK-Prüfung in drei zeitlich getrennten Teilen; Teil~1 nach ca.\ 15 Monaten (20\,\%), Teil~2 zum Ausbildungsende (80\,\%).}

\glossareintrag{Handlungskompetenz}{Bereitschaft und Befähigung, sich sachgerecht, individuell, durchdacht und sozial verantwortlich zu verhalten (Fach-, Sozial- und Selbstkompetenz).}

\glossareintrag{BIBB}{Bundesinstitut für Berufsbildung; gibt Ausbildungsordnungen heraus (www.bibb.de).}

\glossareintrag{IHK}{Industrie- und Handelskammer; zuständige Stelle für IT-Berufe, lässt Betriebe zu und nimmt Prüfungen ab (www.ihk.de).}

\end{document}
